\section{Installation}\label{sec:installation}

PyMORGAN requires \textbf{Python~3.12 or later}. It can be installed directly from \href{https://pypi.org/project/pymorgan/}{PyPI} or installed from source using \href{https://docs.astral.sh/uv/}{\textbf{uv}}.

% ---------------------------------------------------------------------------
\subsection{Quick install via PyPI (Recommended)}\label{sec:install-pypi}

To install the latest official release into your active Python environment:

\begin{lstlisting}[language=bash, style=py]
pip install pymorgan
# or with uv:
uv pip install pymorgan

# Register bundled fonts into matplotlib:
pymorgan-install-fonts
\end{lstlisting}

\noindent You can also launch the graphical interface directly without prior installation using \code{uvx}:

\begin{lstlisting}[language=bash, style=py]
uvx --from pymorgan pymorgan-gui
\end{lstlisting}

% ---------------------------------------------------------------------------
\subsection{Installation from Source}\label{sec:install-source}

The steps below describe the full setup for development or source installation.

\paragraph{Step 1 — Install Python 3.12+.}

\paragraph{Windows.}
Download and run the official installer from
\url{https://www.python.org/downloads/}. Choose \textbf{Python 3.12} or
later, and tick \emph{Add Python to PATH} on the first installer page.

Alternatively, use the \textbf{Windows Store} or
\href{https://github.com/indygreg/python-build-standalone/releases}{Python
Standalone Builds} (recommended for reproducibility).

After the installer finishes, verify Python is on \code{PATH}:

\begin{lstlisting}[language=bash, style=py]
python --version        # should print Python 3.12.x or later
\end{lstlisting}

\paragraph{macOS.}
\begin{lstlisting}[language=bash, style=py]
brew install python@3.12          # Homebrew
# or download from https://www.python.org/downloads/
\end{lstlisting}

\paragraph{Linux (Debian/Ubuntu).}
\begin{lstlisting}[language=bash, style=py]
sudo apt update
sudo apt install python3.12 python3.12-venv python3.12-dev
\end{lstlisting}

% ---------------------------------------------------------------------------
\subsection{Step 2 — Install uv}\label{sec:install-uv}

uv is a single self-contained binary that replaces \code{pip},
\code{venv}, \code{pip-tools} and \code{virtualenv} in one tool. Install
it once, globally:

\paragraph{Windows (PowerShell).}
\begin{lstlisting}[language=bash, style=py]
powershell -ExecutionPolicy ByPass -c "irm https://astral.sh/uv/install.ps1 | iex"
\end{lstlisting}

\paragraph{macOS / Linux.}
\begin{lstlisting}[language=bash, style=py]
curl -LsSf https://astral.sh/uv/install.sh | sh
\end{lstlisting}

\paragraph{Via \code{pip} (any OS).}
\begin{lstlisting}[language=bash, style=py]
pip install uv
\end{lstlisting}

After installing, open a new terminal window and confirm:
\begin{lstlisting}[language=bash, style=py]
uv --version     # e.g. uv 0.6.x
\end{lstlisting}

% ---------------------------------------------------------------------------
\subsection{Step 3 — Obtain the PyMORGAN source}\label{sec:clone}

Clone the repository (or copy the folder to a location of your choice):

\begin{lstlisting}[language=bash, style=py]
git clone https://github.com/RJFernandezTeran/PyMORGAN.git
cd PyMORGAN
\end{lstlisting}

If you do not use git, download and extract the zip archive from GitHub,
then open a terminal inside the extracted folder.

% ---------------------------------------------------------------------------
\subsection{Step 4 — Create the virtual environment and install}\label{sec:venv}

From the repository root:

\begin{lstlisting}[language=bash, style=py]
uv venv                       # creates .venv/ with the system Python (>= 3.12)
uv pip install -e ".[dev]"    # editable install + ruff + pytest
uv run pymorgan-install-fonts # copy bundled fonts into matplotlib
\end{lstlisting}

\noindent\textbf{What each line does:}
\begin{itemize}
  \item \code{uv venv} — creates a \code{.venv/} virtual environment in the
    current directory.  Pass \code{--python 3.12} to pin a specific
    interpreter if multiple versions are installed.
  \item \code{uv pip install -e ".[dev]"} — installs all core dependencies
    (Table~\ref{tab:dependencies}), plus the development extras
    (\code{ruff} linter and \code{pytest} test runner), in \emph{editable}
    mode. Editable mode means changes to the source files take effect
    immediately without reinstalling.
  \item \code{uv run pymorgan-install-fonts} — copies the bundled
    \code{TgTermes} and other fonts used by the plot style profiles into
    matplotlib's font cache. Run this once after a fresh install.
\end{itemize}

\noindent If you only need to run the GUI (no development work):

\begin{lstlisting}[language=bash, style=py]
uv pip install -e .           # core only (no ruff / pytest)
uv run pymorgan-install-fonts
\end{lstlisting}

% ---------------------------------------------------------------------------
\subsection{Step 5 — Verify the installation}\label{sec:verify}

Confirm the package imports and reports its version:

\begin{lstlisting}[language=bash, style=py]
uv run python -c "import pymorgan; print(pymorgan.__version__)"
\end{lstlisting}

Launch the graphical interface:

\begin{lstlisting}[language=bash, style=py]
uv run pymorgan-gui
\end{lstlisting}

Run the full test suite (requires the \code{dev} extras):

\begin{lstlisting}[language=bash, style=py]
uv run pytest        # 112 tests, all should pass
\end{lstlisting}

Run the end-to-end smoke check:

\begin{lstlisting}[language=bash, style=py]
uv run python scripts/run_checks.py
\end{lstlisting}

% ---------------------------------------------------------------------------
\subsection{Dependencies}\label{sec:dependencies}

Table~\ref{tab:dependencies} lists the core runtime dependencies. All are
installed automatically by \code{uv pip install -e .}; no manual
\code{pip install} is needed.

\begin{table}[h]
  \centering
  \begin{tabular}{lll}
    \toprule
    Package & Minimum version & Purpose \\
    \midrule
    \code{numpy}           & 2.5.2  & Array operations \\
    \code{scipy}           & 1.18.0 & Signal processing and fitting \\
    \code{matplotlib}      & 3.11.1 & Plotting back-end \\
    \code{pandas}          & 3.0.5  & Tabular data and I/O \\
    \code{h5py}            & 3.16.0 & HDF5 file I/O (some loaders) \\
    \code{mpl-axes-aligner}& 1.3    & Aligned multi-axis figures \\
    \code{tomlkit}         & 0.15.1 & Settings file read/write \\
    \code{PyQt6}           & 6.11.0 & Graphical interface framework \\
    \code{PyQt6\_sip}      & 13.12.0& Qt bindings interface \\
    \bottomrule
  \end{tabular}
  \caption{Core runtime dependencies (minimum version floors as of August 2026).
    Exact versions are pinned in \code{uv.lock}.}
  \label{tab:dependencies}
\end{table}

\begin{table}[h]
  \centering
  \begin{tabular}{lll}
    \toprule
    Package & Minimum version & Purpose \\
    \midrule
    \code{ruff}      & 0.16.2 & Linter and formatter \\
    \code{pytest}    & 9.1.1  & Test runner \\
    \code{pytest-qt} & 4.5.0  & Qt widget testing helpers \\
    \bottomrule
  \end{tabular}
  \caption{Development extras (\code{uv pip install -e ".[dev]"}).}
  \label{tab:dev-dependencies}
\end{table}

% ---------------------------------------------------------------------------
\subsection{Keeping packages up to date}\label{sec:upgrade}

To upgrade all dependencies to the latest compatible versions:

\begin{lstlisting}[language=bash, style=py]
uv lock --upgrade     # resolve latest versions, update uv.lock
uv sync               # install the resolved versions into .venv
\end{lstlisting}

% ---------------------------------------------------------------------------
\subsection{Console scripts}\label{sec:scripts}

Four console scripts are installed with the package:

\begin{table}[h]
  \centering
  \begin{tabular}{ll}
    \toprule
    Script & Purpose \\
    \midrule
    \code{pymorgan-gui}           & Launch the PyQt6 graphical interface \\
    \code{pymorgan-install-fonts} & Install bundled fonts into matplotlib \\
    \code{pymorgan-edit-gui}      & Open the GUI layout in Qt Designer \\
    \code{pymorgan-settings}      & Standalone settings editor (no data loaded) \\
    \bottomrule
  \end{tabular}
  \caption{Console scripts installed with the package.}
  \label{tab:scripts}
\end{table}

All scripts are run via \code{uv run <script>} (e.g.\
\code{uv run pymorgan-gui}), or directly if the environment is activated.

% ---------------------------------------------------------------------------
\subsection{GUI layout and compiled UI}\label{sec:compiled-ui}

The GUI layout is defined in
\code{src/pymorgan/gui/main\_window.ui} (Qt Designer XML). For performance,
PyMORGAN also ships a pre-compiled Python translation,
\code{main\_window\_ui.py}, generated by \code{pyuic6}. Loading the
compiled module is approximately $1.5\times$ faster than parsing the XML at
runtime.

\textbf{No manual action is needed when editing the \code{.ui} file.} When
\code{main\_window.ui} is saved in Qt Designer and the GUI is next launched,
PyMORGAN detects that the compiled module is stale (via file modification
times) and regenerates it automatically using \code{pyuic6} before loading
the window. If \code{pyuic6} is unavailable or fails, the raw XML is parsed
as a transparent fallback.

To force a manual recompile (e.g.\ for CI):

\begin{lstlisting}[language=bash, style=py]
uv run pyuic6 src/pymorgan/gui/main_window.ui \
              -o src/pymorgan/gui/main_window_ui.py
\end{lstlisting}

% ---------------------------------------------------------------------------
\subsection{Activating the environment (optional)}\label{sec:activate}

All commands in this manual use the \code{uv run} prefix, which
automatically targets the project's \code{.venv} without requiring
activation. If you prefer an activated shell (e.g.\ for interactive work
in a terminal):

\begin{lstlisting}[language=bash, style=py]
# Linux / macOS
source .venv/bin/activate

# Windows (PowerShell)
.venv\Scripts\Activate.ps1

# Windows (Command Prompt)
.venv\Scripts\activate.bat
\end{lstlisting}

Once activated, commands can be run without the \code{uv run} prefix
(e.g.\ \code{pymorgan-gui} instead of \code{uv run pymorgan-gui}).
